% !TeX root = ../navier-stokes-manuscript.tex
% ============================================================
% 3.4 One Endpoint for the Whole Tower
% ============================================================
\subsection{One Endpoint for the Whole Tower}\label{sub:OneEndpointForTheWholeTower}

A restart needs an actual datum at \(T_*\), not a separate endpoint in each
rectangle.  Under \Cref{ass:BodyWholeTerminalRectangleCondition},
\Cref{thm:WholeTerminalCompatibleRectangleEstimate} supplies the complete
spatial row on a hypothetical finite branch.  The equation then supplies
enough time integrability to make the same velocity converge at every finite
height.

\begin{lemma}[The complete spatial row reaches every endpoint topology]\label[lemma]{lem:CompleteSpatialRowEndpoint}
Let \((u,p)\) be the fixed normalized history on \(I_*\).  If
\(u\in\Szero(I_*)\), then for every \(m\in\Nzero\),
\begin{equation}\label{eq:SectionThreeWOneOneEndpoint}
  u_t\in L^1(I_*;H^m),
  \qquad
  u\in W^{1,1}(I_*;H^m)
  \hookrightarrow C([0,T_*];H^m).
\end{equation}
\end{lemma}

\begin{proof}
The pressure-complete transport response can be written without separating
pressure:
\[
  (u\cdot\nabla)u+\nabla p
  =\Leray\nabla\cdot(u\otimes u).
\]
Under the fixed normalization, this is exactly transport plus \(\nabla p\).
Sobolev multiplication and boundedness of the Leray and Riesz multipliers give
\[
  \norm{\Leray\nabla\cdot(u\otimes u)}_{H^m}
  \leq
  C_m\norm{u}_{H^{m+2}}^2.
\]
The complete row, together with \(T_*<\infty\), therefore yields
\begin{align*}
  \norm{\Leray\nabla\cdot(u\otimes u)}_{L^1(I_*;H^m)}
  &\leq
  C_m\norm{u}_{L^2(I_*;H^{m+2})}^{2}
  <\infty,\\
  \norm{\Delta u}_{L^1(I_*;H^m)}
  &\leq
  T_*^{1/2}\norm{u}_{L^2(I_*;H^{m+2})}
  <\infty.
\end{align*}
The momentum equation puts \(u_t\) in \(L^1(I_*;H^m)\).  In particular, for
\(0\leq s<t<T_*\),
\[
  \norm{u(t)-u(s)}_{H^m}
  \leq
  \nu(t-s)^{1/2}\norm{u}_{L^2((s,t);H^{m+2})}
  +C_m\norm{u}_{L^2((s,t);H^{m+2})}^{2}.
\]
Both terms vanish when \(s,t\uparrow T_*\), making the same \(u(t)\) Cauchy
in \(H^m\).  Since the finite interval also embeds \(L^2_tH^m_x\) into
\(L^1_tH^m_x\), \(u\in W^{1,1}(I_*;H^m)\); the Bochner fundamental theorem of
calculus supplies the continuous representative and its endpoint.
\end{proof}

Let \(u_*^{(m)}\) be the endpoint found in \(H^m\).  If \(k>m\), both the
\(H^k\) trace embedded into \(H^m\) and the native \(H^m\) trace represent the
same \(u\) almost everywhere.  Continuity makes them identical through the
closed interval.  Their terminal values are therefore one element
\begin{equation}\label{eq:SectionThreeCommonEndpoint}
  u_*:=u(T_*)
  \in\bigcap_{m=0}^{\infty}H^m(\R^3)
  =\Hinfty(\R^3).
\end{equation}
No uniformity in \(m\) is used.  Divergence passes continuously to the limit,
so \(u_*\) is solenoidal, and the strong \(L^2\) trace preserves finite kinetic
energy.

Apply the same adjacent trace to every temporal coordinate:
\[
  \partial_t^N u\in C([0,T_*];H^m)
  \qquad(N,m\in\Nzero).
\]
For fixed \(N\), spatial compatibility identifies all these traces as
\begin{equation}\label{eq:SectionThreeVelocityJetTraces}
  V_N^*
  :=\lim_{t\uparrow T_*}\partial_t^N u(t)
  \in\Hinfty(\R^3).
\end{equation}
It remains to show that this trace family is the actual Navier--Stokes jet
generated by \(u_*\).

\Needspace{15\baselineskip}
\begin{proposition}[Pressure-complete endpoint jet]\label[proposition]{prop:PressureCompleteEndpointJet}
The endpoint \(u_*\) determines one sequence
\((V_N^*,Q_N^*)_{N\geq0}\subset\Hinfty\times\Hinfty\) by
\begin{align}
  V_0^*&=u_*,\notag\\
  Q_N^*&=R_iR_j\sum_{a=0}^{N}\binom{N}{a}V_{a,i}^*V_{N-a,j}^*,
  \label{eq:SectionThreeEndpointPressureJet}\\
  V_{N+1}^*&=\nu\Delta V_N^*
  -\sum_{a=0}^{N}\binom{N}{a}(V_a^*\cdot\nabla)V_{N-a}^*
  -\nabla Q_N^*.
  \label{eq:SectionThreeEndpointVelocityJet}
\end{align}
This sequence is exactly the family of left traces of the original normalized
velocity and pressure history.
\end{proposition}

\begin{proof}
The space \(\Hinfty=\bigcap_mH^m\) is stable under multiplication and spatial
differentiation, and Riesz transforms preserve every one of its Sobolev
coordinates.  The two recurrences therefore determine a unique next pair at
each finite order.

Fix a finite target height \(m\) and take the trace convergence in a Sobolev
space high enough for every product and derivative in the recurrence.  All
operations are continuous between these finite spaces, so the limit of
\eqref{eq:SectionThreePressureRecurrence}--\eqref{eq:SectionThreeVelocityRecurrence}
is the displayed endpoint recurrence.  Starting from \(V_0^*=u_*\), induction
in \(N\) identifies the recursively generated entries with the left traces.
Since \(m\) was arbitrary, the identity holds in \(\Hinfty\), and spatial
differentiation supplies every mixed endpoint coordinate.
\end{proof}

No new pressure constant appears at the terminal face.  The decaying Riesz
formula fixes every \(Q_N(t)\) on the left, and the limit preserves that
normalization.  Thus inserting \(u_*\) into the normalized recurrence produces
exactly the left mixed jet.
